Post-hoc explanations are increasingly used to audit machine-learning models for spatial bias, and GeoShapley was
designed to separate location effects from other features. We show that a scaffold attack can hide a racial price
penalty from LIME, KernelSHAP and GeoShapley. A home-valuation model for King County, Washington, halves predicted
prices in census tracts where Black and Latino residents exceed a quarter of the population; a detector trained on
each explainer's own queries serves an innocuous model to the explainer and the biased model to real homes. LIME then
shows \FLIMEMatKept\ of the penalty while \FLIMEMatFid\% of real homes still receive it. KernelSHAP and GeoShapley show
\FKernelSHAPMatKept\ and \FGeoShapleyMatKept\ of it at about 90\% fidelity, and the attacker can hide more only by
withholding the penalty from real homes. GeoShapley is no harder to fool than KernelSHAP, and under attack its
location effect shrinks while spurious feature-by-location interactions appear. The difference between explainers is
set by their estimators, not by their treatment of space: efficiency keeps the penalty inside a Shapley explanation,
and the share kept is close to the share of coalition hybrids the detector cannot tell from real homes. \todo{one
clause on generality once the robustness runs are in: model families and audit settings, rule R.} A residual check
that needs no explainer finds the penalised homes (AUC \EsevenAucLagScaffoldTen\ for a 21\% penalty).
